\documentclass[a4paper, 12pt, english]{article}

% \usepackage[portuges]{babel}
\usepackage[utf8]{inputenc}
\usepackage{amsmath,amssymb}
\usepackage{graphicx}
\usepackage{subfig}
\usepackage[colorinlistoftodos]{todonotes}

\usepackage{indentfirst}
\usepackage{verbatim}
\usepackage{textcomp}
\usepackage{gensymb}
\usepackage{float}

\usepackage{relsize}

\usepackage{lipsum}% http://ctan.org/pkg/lipsum
\usepackage{xcolor}% http://ctan.org/pkg/xcolor
\usepackage{xparse}% http://ctan.org/pkg/xparse
\NewDocumentCommand{\myrule}{O{1pt} O{2pt} O{black}}{%
	\par\nobreak % don't break a page here
	\kern\the\prevdepth % don't take into account the depth of the preceding line
	\kern#2 % space before the rule
	{\color{#3}\hrule height #1 width\hsize} % the rule
	\kern#2 % space after the rule
	\nointerlineskip % no additional space after the rule
}
\usepackage[section]{placeins}

\usepackage{booktabs}
\usepackage{colortbl}%
\newcommand{\myrowcolour}{\rowcolor[gray]{0.925}}

\usepackage[obeyspaces]{url}
\usepackage{etoolbox}
\usepackage[colorlinks,citecolor=black,urlcolor=blue,bookmarks=false,hypertexnames=true]{hyperref}

\usepackage{geometry}
\geometry{
	paper=a4paper, % Change to letterpaper for US letter
	inner=3cm, % Inner margin
	outer=3cm, % Outer margin
	bindingoffset=.5cm, % Binding offset
	top=2cm, % Top margin
	bottom=2cm, % Bottom margin
	%showframe, % Uncomment to show how the type block is set on the page
}
\usepackage{fancyhdr}

% Define header and footer
\pagestyle{fancy}
\fancyhf{}
\lhead{ENER 104L}
\rhead{iSciM, Habib University} % Right-aligned page number in the header
\rfoot{\thepage} % Right footer text
%*******************************************************************************%
%************************************START**************************************%
%*******************************************************************************%
\begin{document}

%************************************TITLE PAGE**************************************%
\begin{titlepage}
	\begin{center}
		\textbf{\LARGE Habib University}\\[0.5cm]
		\textbf{\large iSciM}\\[0.2cm]
		\textbf {\large Fall 2023}\\[0.2cm]
		\vspace{20pt}
		\includegraphics[width=5cm]{../habiblogo.jpg}\\[1cm]
		\par
		\vspace{20pt}
		\textbf{\Large ENER 104L RENEWABLE ENERGY}\\
		\vspace{15pt}
		\myrule[1pt][7pt]
		\textbf{\LARGE  LABORATORY REPORT 3}\\
		\vspace{15pt}
		\textbf{\large Calorimetry}\\
		\myrule[1pt][7pt]
		\vspace{25pt}
		\begin{tabular}{@{}p{5cm}p{3cm}@{}}
			\textbf{\large Student Name} & \textbf{\large Student ID} \\
			Ali Asghar Yousuf            & ay06993                    \\ % No1 
			Syed Ibrahim Ali Haider      & sh06565                    \\ % No2
		\end{tabular}

		\vspace{10pt}
		\begin{tabular}{@{}p{5cm}p{3cm}@{}}
			\textbf{\large Group Name} & \textbf{\large Group No.} \\
			Insane Fr                  & 1                         \\
		\end{tabular}

		\vspace{45pt}
		\textbf {\large Lab Instructors:}\\[0.2cm]
		\Large {Paishwa Naqvi}\\[0.1cm]
		\Large {Amber Talat}\\[0.1cm]
	\end{center}

	\par
	\vfill
	\begin{center}
		\textbf{\today}\\
	\end{center}

\end{titlepage}

%************************************TABLE OF CONTENTS**************************************%

%  %Sumário
%  \newpage
%  \tableofcontents
%  \thispagestyle{empty}
%  %End Sumário

%********************************%
%***********SECTION 1************%
%********************************%
\newpage
\section{Objectives}
\begin{itemize}
	\item To determine the potential energy content of different food items compare the
	      energy released by different food items.
	\item To conduct a calorimetry experiment to determine the energy content of food
	      items by capturing the energy released.
	\item To estimate the energy content of food items.
	\item To assess the heat content of food items.
	\item To compare the energy content of different food items.
\end{itemize}

\section{Abstract}
This report explores the energy content of different food items. The experiment
was conducted by burning different food items and capturing the energy released
by them. The energy released was captured by heating water in a calorimeter.
The temperature change of water was used to calculate the energy released by
the food items. The energy released by different food items was compared to
determine the energy content of different food items. The energy content of
different food items was also compared to determine which food items have more
energy content. The experiment was conducted by burning different food items
and capturing the energy released by them. The energy released was captured by
heating water in a calorimeter. The temperature change of water was used to
calculate the energy released by the food items. The energy released by
different food items was compared to determine the energy content of different
food items. The energy content of different food items was also compared to
determine which food items have more energy content.

\section{Result and Analysis}
\begin{table}[H]
	\caption{\label{tab:Table 1} Measurements}
	\centering
	\begin{tabular}{c c c c c c c}
		\toprule
		\textbf{Food}              &
		\textbf{$M_i$ (g)}         &
		\textbf{$M_f$ (g)}         &
		\textbf{$\Delta M$ (g)}    &
		\textbf{$T_i$ (\degree C)} &
		\textbf{$T_f$ (\degree C)} &
		\textbf{$\Delta T$ (\degree C)}                                \\

		\cmidrule[0.4pt](r{0.125em}){1-1}%
		\cmidrule[0.4pt](lr{0.125em}){2-2}%
		\cmidrule[0.4pt](lr{0.125em}){3-3}%
		\cmidrule[0.4pt](lr{0.125em}){4-4}%
		\cmidrule[0.4pt](lr{0.125em}){5-5}%
		\cmidrule[0.4pt](lr{0.125em}){6-6}%
		\cmidrule[0.4pt](l{0.25em}r{0.125em}){7-7}%
		% \midrule
		\textbf{Marshmallow}       & 1.25 & 0.80 & 0.45 & 22 & 32 & 10 \\
		\textbf{Chips}             & 0.50 & 0.09 & 0.41 & 22 & 42 & 20 \\
		\textbf{Cornflakes}        & 0.08 & 0.05 & 0.03 & 24 & 26 & 2  \\
		\textbf{White Beans}       & 0.23 & 0.22 & 0.01 & 24 & 25 & 1  \\
		\textbf{Red Beans}         & 0.37 & 0.34 & 0.03 & 23 & 23 & 0  \\
	\end{tabular}
\end{table}

\begin{table}[H]
	\caption{\label{tab:Table 2} Calculations}
	\centering
	\begin{tabular}{c c c}
		\toprule
		\textbf{Food}              &
		\textbf{$\Delta T$ (\degree C)} &
		\textbf{$\Delta E$ (J)} \\

		\cmidrule[0.4pt](r{0.125em}){1-1}%
		\cmidrule[0.4pt](lr{0.125em}){2-2}%
		\cmidrule[0.4pt](l{0.25em}r{0.125em}){3-3}%
		% \midrule
		\textbf{Marshmallow}       & 10 & 418.4 \\
		\textbf{Chips}             & 20 & 836.8 \\
		\textbf{Cornflakes}        & 2  & 83.68 \\
		\textbf{White Beans}       & 1  & 41.84 \\
		\textbf{Red Beans}         & 0  & 0     \\
	\end{tabular}
\end{table}

\section{Conclusion}
abc

\section{Questions}
\begin{enumerate}
	\item How could you improve your procedure to overcome difference in energy value?

	\item W hat is the name of the process in living organisms that releases energy from
	      food?
	\item What are reactants and products in your experiment, write down in word form?
	\item What are dependent and independent variables in this experiment?
	\item Which food items will release more energy? Are they most rich in fat,
	      carbohydrate, or protein? Which food burns well- protein-rich,
	      carbohydrate-rich, fat rich etc.
	\item Define what a calorie is and compare of low calorie and high calorie food items
	\item Explain why body needs calories and what happens to excess calories?
\end{enumerate}

\end{document}